% year: תשע"ח
% type: מבחן
% semester: ב
% moed: ב
% number: 1ב
% points: 10
% categories: continuity, func-limits
% source: exams/מבחנים/תשעח/מועד ב תשעח סמסטר ב.pdf

\begin{question}
מצאו את כל הנקודות האי רציפות של הפונקציה ותמיינו אותם
\[
f(x)=\frac{1}{25-5^{1/(x-2)}}.
\]
\end{question}

\begin{hint}
הנקודות החשודות הן $x=2$ (שם $\frac1{x-2}$ לא מוגדר) והנקודה שבה המכנה מתאפס. ב-$x=2$ חשבו גבולות חד-צדדיים בנפרד.
\end{hint}

\begin{solution}
\textbf{תחום הגדרה:} צריך $x\neq2$ וגם $5^{1/(x-2)}\neq25$, כלומר $\frac1{x-2}\neq2$, כלומר $x\neq\frac52$. בכל נקודה אחרת $f$ היא הרכבה ומנה של פונקציות רציפות, ולכן רציפה. נקודות אי-הרציפות הן $x=2$ ו-$x=\frac52$.

\textbf{הנקודה $x=2$:}
\begin{itemize}
  \item $x\to2^+$: $\frac1{x-2}\to+\infty$, ולכן $5^{1/(x-2)}\to+\infty$ ו-$f(x)\to0$.
  \item $x\to2^-$: $\frac1{x-2}\to-\infty$, ולכן $5^{1/(x-2)}\to0$ ו-$f(x)\to\frac1{25}$.
\end{itemize}
שני הגבולות החד-צדדיים קיימים, סופיים ושונים: \textbf{אי-רציפות מסוג ראשון (קפיצה)}.

\textbf{הנקודה $x=\frac52$:} המכנה $25-5^{1/(x-2)}$ שואף ל-$0$. עבור $2<x<\frac52$: $\frac1{x-2}>2$, ולכן $5^{1/(x-2)}>25$ והמכנה שלילי: $f(x)\to-\infty$ כאשר $x\to\frac52^-$. עבור $x>\frac52$ המכנה חיובי: $f(x)\to+\infty$ כאשר $x\to\frac52^+$. הגבולות החד-צדדיים אינסופיים: \textbf{אי-רציפות מסוג שני} (ו-$x=\frac52$ אסימפטוטה אנכית).
\end{solution}
